% =====================================================================
% SECTION 4 IS WRITTEN BY THE TEAM.
% Budget: about 2 pages. The other sections take about 4 pages.
% Keep the subsection labels: Sections 3 and 5 refer to them.
% Structure: Setup, Baseline results (the method as written), Improvements.
%   Improvements has one \subsubsection per set of improvements (one per team member).
% Sources: SPEC.md, the runs/ and runs_vanilla/ folders on the cluster,
%          python scripts/aggregate.py runs/   (mean +- std over seeds, next to the paper's numbers)
%          python scripts/trajectory.py runs/<task>   (per-epoch means)
% Writing rules (same as the rest of the report):
%  - One idea per sentence, up to about 25 words. Link sentences so the text flows.
%  - Body text uses only commas and periods. No colons, semicolons or dashes.
%    Write ranges as "14 to 28", not "14-28".
%  - Parentheses only around numbers, citations and technical tags.
%  - Every claim gets a number, and every number points to a table, figure or run folder.
%  - Exact figures as in SPEC.md or aggregate.py. Bold only the milestone results.
%  - Tables carry the numbers. Prose carries the interpretation.
%  - State before/after contrasts explicitly, e.g. "ALL rises from X to Y".
%  - The "paper as written" runs are the baseline (a reproduction of the method as written).
%    Label numbers quoted from the paper. Treat differences under 2 points as noise.
% =====================================================================
\section{Experiments}\label{sec:experiments}

\subsection{Setup}\label{sec:exp-setup}
\TODO{Datasets (Digits M/U/S, Office-31, VisDA) and class splits. ALL and PRIV accuracy. 5 seeds with different one-shot images. GPU and run times.}

\subsection{Baseline results}\label{sec:exp-vanilla}
We first ran a reproduction of the method exactly as the paper describes it, with all four changes of \cref{sec:changes} switched off (\texttt{preset=vanilla}).
Where the paper is silent, the reproduction keeps the choices of \cref{tab:choices}.
\Cref{tab:vanilla} compares it with the numbers reported in the paper on every benchmark.
The reproduction stays below the paper on every task.
The gap in ALL ranges from 14 to 29 points on Digits, 11 to 54 points on Office-31 and about 50 points on VisDA.
PRIV suffers most on Office-31, where it ends between 1.4 and 26.9 against 61.8 to 88.4 in the paper.
The way accuracy is lost differs between the benchmarks, and each failure points to one of our changes.

\begin{table}[t]
  \centering
  \small
  \caption{Reproduction of the method as written (\texttt{preset=vanilla}) against the numbers reported in the paper. ALL and PRIV accuracy in \% at the last epoch, mean over 5 seeds. Office 10/20 is the setting $|\Cs|=10, |\Ct|=20$, and Office 20/31 is $|\Cs|=20, |\Ct|=31$. Runs are in \texttt{runs\_vanilla/}.}
  \label{tab:vanilla}
  \begin{tabular}{@{}llcccc@{}}
    \toprule
    & & \multicolumn{2}{c}{Reproduction} & \multicolumn{2}{c}{Paper (reported)} \\
    \cmidrule(lr){3-4}\cmidrule(lr){5-6}
    Setting & Task & ALL & PRIV & ALL & PRIV \\
    \midrule
    Digits & M$\to$U & 43.7 & 50.5 & 72.5 & 71.6 \\
           & U$\to$M & 47.2 & 60.7 & 69.4 & 68.6 \\
           & S$\to$M & 32.8 & 46.5 & 46.4 & 45.7 \\
    \midrule
    Office 10/20 & A$\to$D & 45.2 & 16.9 & 72.2 & 72.6 \\
                 & A$\to$W & 51.3 & 26.9 & 62.1 & 62.0 \\
                 & D$\to$A & 25.2 & 3.2  & 62.6 & 61.8 \\
                 & D$\to$W & 56.1 & 17.8 & 78.5 & 78.7 \\
                 & W$\to$A & 20.9 & 7.4  & 62.1 & 62.4 \\
                 & W$\to$D & 54.3 & 19.6 & 87.8 & 87.6 \\
    \midrule
    Office 20/31 & A$\to$D & 47.7 & 16.3 & 73.3 & 73.1 \\
                 & A$\to$W & 42.5 & 20.9 & 63.6 & 62.6 \\
                 & D$\to$A & 9.9  & 4.4  & 64.1 & 64.3 \\
                 & D$\to$W & 53.2 & 17.7 & 80.3 & 79.4 \\
                 & W$\to$A & 14.7 & 1.4  & 63.7 & 62.8 \\
                 & W$\to$D & 60.0 & 12.7 & 89.5 & 88.4 \\
    \midrule
    VisDA & Syn$\to$Real & 6.2 & 16.1 & 56.4 & 56.3 \\
          & Real$\to$Syn & 7.7 & 13.9 & 55.8 & 55.7 \\
    \bottomrule
  \end{tabular}
\end{table}

\paragraph{Digits.}
On Digits, shared samples move into the private classes.
On M$\to$U, accuracy on the shared classes is already 47.2 before any adaptation loss is applied.
At that point the only differences from the source model are the pretrained projections and the source BatchNorm statistics.
During the first epoch, the nearest-guide pseudo-labels are only 37.5 to 43.9\% accurate.
Without target BatchNorm statistics, shared target samples lie closer to the one-shot private guides than to the shared guides.
$\mathcal{L}_{a1}$ then trains $f_t$ on these labels for every target sample.
By epoch 30, 57 to 60\% of the shared samples are predicted as private, and shared accuracy falls to 20.0 on S$\to$M.
PRIV stays at 46.5 to 60.7 only because most samples are predicted as private.

\paragraph{Office-31.}
On Office-31 the failure runs the other way, and private samples move into the shared classes.
PRIV peaks within the first 10 epochs and then decays steadily.
On A$\to$D it peaks at 35.6 in epoch 10 and ends at 16.9, and on D$\to$A it falls from 36.9 in epoch 2 to 3.2.
At epoch 60, 53 to 94\% of the private samples are predicted as shared in the 10/20 setting.
$\mathcal{L}_{a1}$ fine-tunes the whole ResNet-50 on noisy labels for every target sample, so the features drift towards the shared guides.
With Amazon as the target, shared accuracy collapses as well, which brings ALL on D$\to$A in the 20/31 setting down to 9.9.
Even the best epoch, chosen with target labels that the setting does not allow, stays below the paper.
On W$\to$D in the 10/20 setting, ALL peaks at 72.6 against 87.8 in the paper.

\paragraph{VisDA.}
On VisDA the model collapses onto a single class in both directions.
Class-averaged ALL is 8.5 on Syn$\to$Real and 8.4 on Real$\to$Syn, close to the 8.3 of a constant prediction over 12 classes.
The pipeline itself works, since the source model reaches 37.3 on the shared real images before adaptation, against 16.7 for chance.
The nearest-guide pseudo-labels, however, are only 10.4\% accurate at that point.
The shared guides are built from synthetic features, while the private guides come from real one-shot images.
In seed 0, the pretrained projections also remain 24\% away from the identity, so the shared guides are misplaced from the start.
By epoch 10, shared accuracy on Syn$\to$Real is 0.0, and every target image is assigned to one private class.

\paragraph{Summary.}
The three failures share one mechanism.
Wrong pseudo-labels in the first epochs are reinforced by $\mathcal{L}_{a1}$, which trains on every target sample.
Source BatchNorm statistics and imperfect projections make these early labels worse, and fine-tuning the whole backbone lets the errors reshape the features.
Each of these causes corresponds to one change in \cref{sec:changes}.

\subsection{Improvements}\label{sec:exp-improvements}
Each set of improvements below is compared with the baseline of \cref{sec:exp-vanilla} and with the numbers reported in the paper.
% Each team member adds their own improvements as a \subsubsection below,
% with its own label (sec:exp-imp-<topic>), a list of the changes and a results table.

% ---------------------------------------------------------------------
% Improvements by Nikita: the four changes of Section 3 (sec:changes).
% Empty cells in tab:imp-pseudo are runs that have not finished yet.
% ---------------------------------------------------------------------
% TEMPORARY title naming the author. Final title: "Stabilising the pseudo-labels".
\subsubsection{Nikita's improvements}\label{sec:exp-imp-pseudo}
This set combines the four changes of \cref{sec:changes}.
Each one targets a cause of the baseline failures, and all four are on by default in our configurations.
\begin{itemize}
  \item \textbf{Target BatchNorm statistics.} Before adaptation, AdaBN re-estimates the BatchNorm statistics of $f_t$ on the unlabelled target set (\texttt{adapt.adabn}).
  \item \textbf{Identity projections.} $f_e$ and $f_d$ take the residual form $\uu+h(\uu)$ with a linear last layer initialised to zero, so both start as the exact identity (\texttt{model.proj\_residual}).
  \item \textbf{$\mathcal{L}_{a1}$ on the confident set.} The $\mathcal{L}_{a1}$ batches come from $\Bt$ instead of the whole target set (\texttt{adapt.la1\_samples}).
  \item \textbf{Frozen ResNet-50.} On Office-31 and VisDA the ResNet-50 keeps its ImageNet weights in both stages, while its BatchNorm statistics still adapt (\texttt{backbone\_lr\_mult}). Digits uses a LeNet trained from scratch, so this change does not apply there.
\end{itemize}
\Cref{tab:imp-pseudo} reports the results next to the baseline and the paper.
Runs are in \texttt{runs/}, and empty cells are runs that have not finished yet.

\begin{table}[t]
  \centering
  \footnotesize
  \setlength{\tabcolsep}{3.5pt}
  \caption{Results with the four changes of \cref{sec:changes} against the baseline of \cref{tab:vanilla} and the numbers reported in the paper. ALL and PRIV accuracy in \% at the last epoch, mean over 5 seeds. Empty cells are runs that have not finished yet. Bold marks tasks where ALL matches or exceeds the paper.}
  \label{tab:imp-pseudo}
  \begin{tabular}{@{}llcccccc@{}}
    \toprule
    & & \multicolumn{2}{c}{Improved} & \multicolumn{2}{c}{Baseline} & \multicolumn{2}{c}{Paper (reported)} \\
    \cmidrule(lr){3-4}\cmidrule(lr){5-6}\cmidrule(lr){7-8}
    Setting & Task & ALL & PRIV & ALL & PRIV & ALL & PRIV \\
    \midrule
    Digits & M$\to$U & \textbf{71.9} & 55.3 & 43.7 & 50.5 & 72.5 & 71.6 \\
           & U$\to$M & 65.5 & 51.0 & 47.2 & 60.7 & 69.4 & 68.6 \\
           & S$\to$M & \textbf{58.8} & 44.0 & 32.8 & 46.5 & 46.4 & 45.7 \\
    \midrule
    Office 10/20 & A$\to$D & \textbf{70.3} & 48.5 & 45.2 & 16.9 & 72.2 & 72.6 \\
                 & A$\to$W & \textbf{67.1} & 44.4 & 51.3 & 26.9 & 62.1 & 62.0 \\
                 & D$\to$A & 41.2 & 11.4 & 25.2 & 3.2  & 62.6 & 61.8 \\
                 & D$\to$W & 72.3 & 48.6 & 56.1 & 17.8 & 78.5 & 78.7 \\
                 & W$\to$A & 43.4 & 12.2 & 20.9 & 7.4  & 62.1 & 62.4 \\
                 & W$\to$D & 73.5 & 52.1 & 54.3 & 19.6 & 87.8 & 87.6 \\
    \midrule
    Office 20/31 & A$\to$D & 65.3 & 26.6 & 47.7 & 16.3 & 73.3 & 73.1 \\
                 & A$\to$W & \textbf{66.6} & 25.5 & 42.5 & 20.9 & 63.6 & 62.6 \\
                 & D$\to$A &      &      & 9.9  & 4.4  & 64.1 & 64.3 \\
                 & D$\to$W &      &      & 53.2 & 17.7 & 80.3 & 79.4 \\
                 & W$\to$A &      &      & 14.7 & 1.4  & 63.7 & 62.8 \\
                 & W$\to$D &      &      & 60.0 & 12.7 & 89.5 & 88.4 \\
    \midrule
    VisDA & Syn$\to$Real &  &  & 6.2 & 16.1 & 56.4 & 56.3 \\
          & Real$\to$Syn &  &  & 7.7 & 13.9 & 55.8 & 55.7 \\
    \bottomrule
  \end{tabular}
\end{table}

\paragraph{Digits.}
ALL rises by 18 to 28 points over the baseline on all three tasks.
It matches the paper on M$\to$U (71.9 against 72.5) and exceeds it on S$\to$M (58.8 against 46.4).
PRIV is lower than the baseline on U$\to$M and S$\to$M, but the baseline PRIV was inflated.
The baseline predicted most target samples as private, which raises PRIV at the cost of shared accuracy.
Moving only $\mathcal{L}_{a1}$ to $\Bt$ raises shared accuracy on M$\to$U at epoch 30 from 52.8 to 83.4.

\paragraph{Office-31.}
ALL rises by 16 to 25 points over the baseline on every finished transfer.
It matches the paper on A$\to$D in the 10/20 setting (70.3 against 72.2) and exceeds it on A$\to$W in both settings.
Shared accuracy stays at 77 to 100 throughout adaptation, whereas the baseline fell to 36 to 50 with Amazon as the target.
PRIV also improves on every transfer, but it remains the main gap at 11.4 to 52.1 against 61.8 to 87.6 in the paper.
The remaining loss is private samples predicted as shared, which is 16 to 70\% of them at epoch 30.
PRIV peaks between epochs 8 and 30 and then decays slowly, for example from 58.3 to 52.1 on W$\to$D.
Amazon targets are the weakest, with PRIV 11.4 on D$\to$A and 12.2 on W$\to$A.
There, only 54 to 60\% of the private part of $\Bt$ is correct, against 73 to 93\% with DSLR or Webcam as the target.
The 20/31 setting gives lower PRIV than the 10/20 setting on the same transfer, 26.6 against 48.5 on A$\to$D.
With 20 shared classes, more shared guides compete for each private sample, and 62 to 65\% of private samples are predicted as shared at epoch 30.

\paragraph{VisDA.}
We have not yet run these changes on VisDA.
